\section{Introduction}
\label{sec:introduction}

Vietnamese social-media language is a consequential setting for affective NLP: users express evaluations in short, informal, context-dependent text, while platforms and downstream analyses often reduce those expressions to surface polarity. That reduction is useful for ordinary sentiment tasks, but it is insufficient when the communicative force of an utterance depends on implication, incongruity, figurative language, language alternation, or ridicule. Vietnamese pretrained and social-media-oriented models provide valuable modelling context \citep{nguyen-tuan-nguyen-2020-phobert,conneau-etal-2020-unsupervised,nguyen-etal-2023-visobert}, and Vietnamese sentiment and emotion resources have enabled important evaluations \citep{nguyen-etal-2018-uit-vsfc,ho-etal-2020-uit-vsmec,tran-etal-2026-vigoemotions}. However, their task definitions do not jointly evaluate the pragmatic phenomena considered here.

Pragmatic sentiment is not a single extension of polarity. Implicit sentiment can convey evaluation without an overt sentiment word; sarcasm and irony can reverse or complicate literal cues; idioms and other figurative expressions can require culturally grounded interpretation; code-switching can alter the lexical and discourse signals available to a classifier; and mocking can signal stance through ridicule rather than an explicit polarity marker. Work on verbal irony and Vietnamese figurative-language evaluation reinforces the importance of context and phenomenon-specific coverage \citep{ghosh-etal-2019-verbal-irony,rem-lab-2026-vivid}. These phenomena can co-occur and remain difficult to capture with a generic sentiment benchmark. At the same time, this observation is not a claim that prior Vietnamese resources are inadequate for their intended purposes, nor that the six labels exhaust Vietnamese pragmatics.

The empirical gap is therefore specific: a provenance-bounded evaluation is needed that brings these six phenomena into one fixed Vietnamese setting, distinguishes them from auxiliary polarity and emotion labels, and compares several realistic model families under the same protocol. ViPragSent supplies 12,000 adjudicated records, with an 8,000/2,000/2,000 train/development/test split. Its gold corpus comprises 10,000 retained local ViSoBERT-export records and 2,000 VIVID-derived, context-augmented candidate records; source text remains private and the latter provenance retains authorization and licence limitations. The benchmark is consequently an evaluation resource with explicit access boundaries, not a public claim to redistribute all underlying text.

The comparison also needs a distinction between model family and evaluation result. Auxiliary-task learning can jointly model related linguistic signals with targeted sentiment, but its benefit is an empirical question rather than a default expectation \citep{moore-barnes-2021-multi}. Likewise, prompted and instruction-tuned models require task- and language-specific evaluation rather than assumed superiority \citep{toukmaji-flanigan-2025-adapting,dou-etal-2024-sailor}. We therefore retain complete best-tuned baselines from PhoBERT, XLM-R-large, Sailor-7B, and Vistral-7B, and assess a frozen label-wise ViPragSent-Final ensemble. The final system is not offered as a novel ensemble architecture: it is a development-selected composition designed to test whether different retained sources and thresholds yield a stronger aggregate result while keeping label-level trade-offs visible.

The selection boundary is central. ViPragSent-Final uses five-fold development out-of-fold evidence to select its label-wise sources and thresholds, freezes the configuration, and then evaluates it once on the canonical test partition. This keeps development selection distinct from test evaluation, while the stored predictions, configuration, and artifact hashes support artifact-level traceability rather than an independent training rerun. Historical 73.7469 ViPragSent configurations, ablations, and diagnostics remain useful versioned evidence, but they are not interchangeable with the promoted final system. Similarly, the stopped true-anchor arbiter cycle did not access canonical test labels and is not promoted.

This study addresses two questions:
\begin{description}
  \item[\textbf{RQ1.}] How does ViPragSent-Final compare overall and label by label with the complete retained best-tuned baselines on the fixed pragmatic evaluation?
  \item[\textbf{RQ2.}] What does the final system's mixed pattern of label-wise leadership reveal about interpreting an aggregate macro-F1 gain?
\end{description}

The answers require both an aggregate comparison and transparent exceptions. ViPragSent-Final reaches 84.3883 macro pragmatic F1 on the post-freeze canonical test evaluation. Vistral-7B, the strongest complete best-tuned baseline, reaches 82.8250, yielding an aggregate difference of +1.5633 points. This result does not establish universal superiority: the final system remains below the strongest retained baseline result for irony, idiom/figurative language, and code-switching. RQ2 therefore treats aggregate performance as a summary to be interpreted alongside label-specific leadership, not as a substitute for it.

The paper makes five contributions:
\begin{enumerate}
  \item It defines a six-label Vietnamese pragmatic-sentiment evaluation with separate polarity and emotion auxiliaries, fixed splits, and a private-data governance boundary.
  \item It documents the annotation, adjudication, provenance, and access limitations needed to interpret the 12,000-record gold corpus.
  \item It provides a common comparative evaluation across Vietnamese, multilingual, social-media-oriented, 7B instruction-tuned, and prompted baseline families, with explicit seed/run qualifications.
  \item It reports ViPragSent-Final as a frozen, development-selected label-wise system evaluated once after freeze, rather than as an unqualified architectural advance.
  \item It reports both the +1.5633 aggregate advantage and all label-level gaps, linking the resulting claims to predictions, configuration, and validation artifacts.
\end{enumerate}

Section~\ref{sec:related-work} positions these choices against related resources and model families. Sections~\ref{sec:task-data}--\ref{sec:results} specify the task, protocol, and comparison; the analysis then separates final-system conclusions from historical and exploratory diagnostics.
